% NSF MSPRF Facilities, Equipment and Other Resources (NSF 23-603) -- DRAFT
%
% Research.gov requires this section in all postdoctoral fellowship proposals (a file reading
% "Not applicable" would pass the automated check). A short description is better here: the
% MSPRF review criteria include "suitability and availability of the sponsoring scientist and
% other colleagues, as well as other conditions at the proposed host institution such as
% adequate space, basic services, and supplies." PAPPG II.D.2.g: narrative only, describe only
% directly applicable resources, no quantifiable financial information. No page limit.
%
% The statements about MIT come from Maulik's sponsoring scientist statement or the Project
% Description; the office was confirmed by the PI (2026-10-20), who needs no computing.

\documentclass[11pt]{article}

\usepackage[margin=1in]{geometry}
\usepackage{amsmath, amssymb}
\usepackage{microtype}
\usepackage{enumitem}
\usepackage[compact]{titlesec}

\titleformat{\section}{\large\bfseries}{\thesection.}{0.5em}{}
\titlespacing*{\section}{0pt}{1.2ex plus .4ex minus .2ex}{0.5ex plus .2ex}
\setlist{nosep, leftmargin=*}

% Block paragraphs, as in the Project Description.
\setlength{\parindent}{0pt}
\setlength{\parskip}{0.3\baselineskip plus 2pt minus 1pt}

\pagestyle{empty}

\begin{document}

\begin{center}
  {\Large\bfseries Facilities, Equipment and Other Resources}
\end{center}

The proposed research is in pure mathematics and requires no laboratory facilities or
specialized equipment. Its needs are office space, access to the mathematical literature,
and an active research community, all of which are available at the host institution.

\section*{Host institution}

The fellowship will be held in the Department of Mathematics at the Massachusetts Institute
of Technology (MIT), with Professor Davesh Maulik as sponsoring scientist. Professor Maulik
will be in residence throughout the fellowship and will meet with the PI weekly. The PI will
have an office in the department.

\section*{Library and computing}

The MIT Libraries provide access to the research literature in mathematics, including
MathSciNet. The project needs no computing resources beyond a personal laptop.

\section*{Scientific environment}

MIT has a large and active community in algebraic geometry and neighboring fields, including
moduli theory, enumerative geometry, and geometric representation theory. The PI will take
part in the department's algebraic geometry seminar and preprint seminar, which Professor
Maulik has invited the PI to help organize. The greater Boston area adds several geometers
with Hodge-theoretic interests, notably Mihnea Popa at Harvard University, and many seminars
in algebraic geometry and related fields.

\section*{Professional development and outreach}

The department offers an intensive teacher-training program for new postdoctoral
researchers, as well as workshops on grant writing and undergraduate advising, in which the
PI will participate. The Broader Impacts activities described in the Project Description
will draw on local programs, including the MIT Summer Research Program, the Harvard--MIT
Mathematics Tournament (HMMT), and Boston-area math circles.

\end{document}
